\chapter{Implementace}
\label{ch:implementation}
Tato kapitola dokumentuje implementaci systému.
Systém sestává ze~čtyř samostatně nasazovaných komponent.
Jádrem je serverový backend, jenž obsluhuje veškerou doménovou logiku včetně integrací s~externími systémy státní správy a~komerčních poskytovatelů.
Klientská strana je realizována jednotnou aplikací rozdělenou na 4 části, jak je popsáno v \ref{sec:komponenty}.
Dvě pomocné Windows služby zajišťují řízení vjezdových závor a~obsluhu lokálního tisku.
Následující sekce popisují tyto komponenty v~uvedeném pořadí.

\section{Backend a~datová vrstva}
\label{sec:impl-backend}
Serverová část je realizována na~platformě \texttt{.NET~10}\footnote{.NET -- dostupné z:~\url{https://dotnet.microsoft.com/}.} v jazyku \texttt{C\#}; verze balíčků jsou centralizovány mechanizmem centrální správy.
Z~ekosystému třetích stran jsou nejvýznamnější knihovny \texttt{Scrutor}\footnote{Scrutor -- dostupné z:~\url{https://github.com/khellang/Scrutor}.} (registrace a~dekorování služeb pro \texttt{DI}), \texttt{FluentValidation}\footnote{FluentValidation -- dostupné z:~\url{https://fluentvalidation.net/}.} (validace), \texttt{Serilog}\footnote{Serilog -- dostupné z:~\url{https://serilog.net/}.} s~výstupem do~agregátoru \texttt{Seq}\footnote{Seq -- dostupné z:~\url{https://datalust.co/seq}.} (logování), \texttt{Npgsql.EFCore.PostgreSQL}\footnote{Npgsql Entity Framework Core Provider -- dostupné z:~\url{https://www.npgsql.org/efcore/}.} (databáze), \texttt{RazorLight}\footnote{RazorLight -- dostupné z:~\url{https://github.com/toddams/RazorLight}.} a~\texttt{Playwright}\footnote{Playwright for .NET -- dostupné z:~\url{https://playwright.dev/dotnet/}.} (generování PDF), \texttt{MailKit}\footnote{MailKit -- dostupné z:~\url{https://github.com/jstedfast/MailKit}.} (transakční elektronická pošta) a~\texttt{QRCoder}\footnote{QRCoder -- dostupné z:~\url{https://github.com/codebude/QRCoder}.} (QR kódy).
Analytický režim je nastaven na~nejpřísnější úroveň a~varování statického analyzátoru \texttt{SonarAnalyzer.CSharp}\footnote{SonarAnalyzer for C\# -- dostupné z:~\url{https://github.com/SonarSource/sonar-dotnet}.} jsou eskalována na~chyby přepínačem \texttt{TreatWarningsAsErrors}.

Projekt je strukturován do~čtyř sestavení odpovídajících vrstvám čisté architektury: doménového jádra (\texttt{Domain}) bez~frameworkových závislostí, aplikační vrstvy s~obsluhami příkazů a~dotazů a~porty pro externí závislosti, infrastrukturní vrstvy implementující perzistenci nad \texttt{EF~Core} a~integrace s~externími systémy, a~prezentační vrstvy (\texttt{Web.Api}) vystavující HTTP rozhraní.
Kompozice mezi vrstvami je řešena pomocí \texttt{DI} podle principu obrácení závislostí.
Toto dělení je vynuceno sadou architekturních testů postavených nad knihovnou \texttt{NetArchTest}\footnote{NetArchTest -- dostupné z:~\url{https://github.com/BenMorris/NetArchTest}.}, jež ověřují, že doménová vrstva nezávisí na~žádné jiné a~že vrstva infrastrukturní není referencována z~prezentační.

Pipeline middlewaru je nakonfigurována v~tomto pořadí: kontext pro strukturované logování, globální obsluha výjimek prostřednictvím rozhraní \texttt{IExceptionHandler}, vyhodnocení CORS, autentizace a~autorizace.
Globální handler převádí nezachycené výjimky na~dokument \emph{Problem Details} ve~formátu \texttt{RFC~9457}.
Endpointy jsou registrovány pod prefixem \texttt{/api}.
Specifikace API je generována ve~formátu \emph{OpenAPI}\,3\footnote{OpenAPI Specification -- dostupné z:~\url{https://www.openapis.org/}.} a~slouží jako referenční kontrakt pro vývoj klientských aplikací; prohlížení je možné pomocí interaktivního prohlížeče API \texttt{Scalar}\footnote{Scalar -- dostupné z:~\url{https://scalar.com/}.}.

Pro explicitní organizaci endpointů byl zaveden vlastní vzor \texttt{IEndpoint}, sdružující cestu, HTTP metodu, autorizační požadavek a~obsluhu do~jedné jednotky.
Při startu aplikace je odpovídající sestavení prohledáno reflexí, všechny nalezené implementace jsou zaregistrovány do~kontejneru závislostí a~jejich registrační metoda je zavolána nad routovacím systémem.

Z~endpointu řízení přechází do~aplikační vrstvy.
Ta nevyužívá běžně užívanou knihovnu \texttt{MediatR}\footnote{MediatR -- dostupné z:~\url{https://github.com/jbogard/MediatR}.}, ale vlastní abstrakce \texttt{ICommand}, \texttt{IQuery} a~odpovídající rozhraní handlerů, jejichž návratovou hodnotou je typ \texttt{Result} nesoucí buď úspěšný výsledek operace, nebo strukturovanou chybu.
Handlery jsou registrovány automaticky knihovnou \texttt{Scrutor} a~obaleny dvěma průchozími dekorátory.
Logovací dekorátor (vnější) zaznamenává začátek a~výsledek operace s~kontextem typu příkazu.
Validační dekorátor (vnitřní) spouští před~vlastním handlerem všechny registrované validátory \texttt{FluentValidation} dohledané k~danému typu příkazu nebo dotazu; při neúspěchu vrací typ \texttt{Result} s~chybou kategorie \texttt{Validation}, aniž by byl handler vůbec zavolán.
Validátory jsou umístěny v~témže adresáři jako odpovídající příkaz.

Zpracování chyb je dvouvrstvé.
Očekávané chybové stavy (chybějící entita, porušení obchodního invariantu, selhání validace) jsou v~handleru reprezentovány návratovou hodnotou \texttt{Result.Failure}.
Endpoint ji mapuje pomocnou metodou \texttt{CustomResults.Problem} na~odpovídající stavový kód a~dokument \emph{Problem Details} podle kategorie chyby: \texttt{Validation} a~\texttt{Problem} produkují 400, \texttt{NotFound} produkuje 404 a~\texttt{Conflict} produkuje 409.

Neočekávané výjimky procházejí globálním exception handlerem, který zvlášť rozeznává \texttt{BadHttpRequestException} (mapovanou na~400).
Výjimky vyvolané perzistentní vrstvou, \texttt{DatabaseConstraintViolationException}, jsou mapovány na stavový kód~409 s~uvedením názvu porušeného omezení.
Veškeré ostatní výjimky jsou logovány v~úrovni \texttt{Error} a~vracejí anonymizovanou odpověď 500.

Autentizace stojí na~knihovně \texttt{ASP.NET~Core~Identity}\footnote{ASP.NET Core Identity -- dostupné z:~\url{https://learn.microsoft.com/en-us/aspnet/core/security/authentication/identity}.}, jejíž aktuální vydání přímo zpřístupňuje úložiště pro bezheslové autentizátory podle standardu \emph{WebAuthn}, viz \ref{sec:webauthn}, čímž odpadá potřeba externí knihovny.
Po~úspěšném ověření vystavuje aplikace pár neprůhledných (opaque) bearer tokenů schématem \texttt{AddBearerToken}\,--\,token přístupový a~obnovovací.
Na~rozdíl od~formátu JWT není jejich obsah pro klienta čitelný a~server jej rekonstruuje deserializací chráněnou mechanizmem \emph{Data~Protection}\footnote{ASP.NET Core Data Protection -- dostupné z:~\url{https://learn.microsoft.com/en-us/aspnet/core/security/data-protection/introduction}.}.
Doba platnosti tokenů je konfigurovatelná pomocí proměnného prostředí.
Pro vývojový režim je předvídatelné autentizační schéma \texttt{NoAuthHandler} aktivovatelné argumentem \texttt{-{}-no-auth}; jeho použití v~produkčním prostředí je zakázáno výjimkou při startu.
Autorizace je role-based, role jsou definovány konstantami ve~třídě \texttt{Roles} a~na~jednotlivé endpointy či skupiny cest se~aplikují metodou \texttt{HasRole}.

Perzistence je řešena přes \texttt{Entity~Framework~Core}\footnote{Entity Framework Core -- dostupné z:~\url{https://learn.microsoft.com/en-us/ef/core/}.} s~ovladačem \texttt{Npgsql}\footnote{Npgsql -- dostupné z:~\url{https://www.npgsql.org/}.} nad databází \texttt{PostgreSQL}\footnote{PostgreSQL -- dostupné z:~\url{https://www.postgresql.org/}.} a~schéma je verzováno migracemi nástroje \texttt{dotnet~ef}.
Doménový model využívá vlastněné typy (\emph{owned types}) pro hodnotové objekty jako adresy nebo datové intervaly, čímž je zachována zapouzdřenost beze ztráty relační reprezentace v~podobě sloupců cílové tabulky.
\texttt{ApplicationDbContext} je registrován s~rozsahem požadavku (\texttt{Scoped}), aby každý souběžně zpracovávaný HTTP požadavek dostal vlastní instanci a~vyhnul se~sdílení tracking stavu mezi vlákny, jež \texttt{EF~Core} nepodporuje.

Tato třída přetěžuje metodu \texttt{SaveChangesAsync}.
Před uložením označuje časová razítka u~entit hostů a~shromažďuje doménové události evidované entitami v~objektu \texttt{ChangeTracker}.
Po~úspěšném zápisu jsou tyto události synchronně předány dispatcheru v~rámci téhož DI scope; ten reflexí dohledá odpovídající handlery a~vykoná je.
Volba synchronního dispatche bez~vzoru \emph{outbox} je vědomá: handlery doménových událostí v~tomto systému neprovádějí žádnou meziprocesovou komunikaci a~jejich případné selhání je akceptovatelné jako stav vyžadující manuální kompenzaci.

Explicitní transakce přes \texttt{BeginTransactionAsync} jsou použity pouze v~generátorech monotónně rostoucích čísel a tvoreb rezervací.
Souběžnost zápisů je jinak řešena na~úrovni databáze prostřednictvím unikátních omezení; kolize na~obchodních identifikátorech tak vyústí v~chybu, kterou \texttt{SaveChangesAsync} odchytí a~přemapuje na~typovanou výjimku popsanou výše.
Aktuální čas je v~aplikační i~doménové vrstvě výhradně získáván třídou implementující rozhraní \texttt{IDateTimeProvider}; přímé použití \texttt{DateTime.UtcNow} je v~těchto vrstvách zakázáno architekturním testem skenujícím IL kód z důvodu testování.

Strukturované logování je zajištěno knihovnou \texttt{Serilog}, jejíž konfigurace je čtena přímo z~aplikační konfigurace a~výstup je ve výchozím nastavení směrován do~agregátoru \texttt{Seq}.
Každý požadavek je obohacen vlastnostmi z~\texttt{RequestContextLoggingMiddleware} a~standardním záznamem \texttt{UseSerilogRequestLogging}, takže korelace události s~konkrétním požadavkem je v~agregátoru triviální.
Logovací dekorátor v~aplikační vrstvě navíc vkládá do~\texttt{LogContext} informaci o~chybě vrácené handlerem; každý záznam o~selhání obchodní operace je tak propojen s~vyvolávajícím požadavkem.

Konfigurace prostředí je celá vázána přes mechanizmus \texttt{IOptions<T>} a~jednotně validována při startu aplikace.
Chybějící povinné údaje způsobí selhání aplikace již při startu, aniž by mohla obsloužit první požadavek.
Typy options jsou definovány pro každou ohraničenou oblast samostatně (nastavení kempu, retenční politiky, frontend, SMTP, tokeny, externí integrace).
Citlivé hodnoty (hesla, klíče, hesla certifikátu) jsou očekávány v~konfiguračních zdrojích prostředí, nikoli ve~zdrojovém kódu.

Dokumenty (vyúčtování, identifikační nálepky pro hosty, finanční uzávěrky) jsou generovány dvoufázově.
Razor šablony zabudované do~projektu infrastruktury jako embedded resources jsou nejprve renderovány knihovnou \texttt{RazorLight} s~podporou lokalizace.
Výsledné HTML je následně převedeno na~PDF knihovnou \texttt{Microsoft.Playwright} v~režimu headless Chromium.
Instance \texttt{IPlaywright} a~\texttt{IBrowser} jsou držené jako singletony po~dobu života procesu, aby se~předešlo opakovanému startu prohlížeče při každém požadavku na~PDF.
QR kódy potřebné pro některé dokumenty jsou vykreslovány knihovnou \texttt{QRCoder}.

\section{Integrace externích systémů}
\label{sec:impl-external}
Integrace s~externími systémy jsou kompoziční součástí infrastrukturní vrstvy popsané v~předchozí sekci; aplikační vrstva s~nimi pracuje výhradně přes doménově orientovaná rozhraní.
Komunikace s~protistranami je v~drtivé většině realizována klientem \texttt{HttpClient} registrovaným jako pojmenovaná konfigurace v~továrně \texttt{IHttpClientFactory}.
Tento mechanizmus centralizuje konfiguraci časových limitů, výchozích adres a~middlewarů a~zaručuje správné životní cykly podkladových socketů.
Každá externí služba je obalena vlastní třídou implementující doménově orientované rozhraní, takže aplikační vrstva s~konkrétní protistranou nepracuje.
Veškeré konfigurační hodnoty (adresy, přihlašovací údaje, cesty k~certifikátům) jsou vázány s~validací při startu aplikace.

Dotazy do~registru ekonomických subjektů \texttt{ARES} vracejí JSON odpovědi deserializované knihovnou \texttt{System.Text.Json}.
Veřejná část registru \texttt{RÚIAN} je publikována ČÚZK jako služba \texttt{ArcGIS~REST} a~je oslovována vlastní třídou \texttt{RuianAddressSuggester}.
Adresní našeptávač portálu \texttt{Mapy.cz} (vyžadující API klíč) je zapojen do~řetězce za~\texttt{RÚIAN} jako záloha pro adresy, jež primární zdroj nepokrývá.
Výsledky obou našeptávačů jsou kvůli zatížení protistran kešovány do~databáze \texttt{Redis}\footnote{Redis -- dostupné z:~\url{https://redis.io/}.}.

Volání služby \texttt{eDoklady} pro zpřístupnění digitálního stejnopisu občanského průkazu vyžaduje vzájemnou TLS autentizaci kvalifikovaným serverovým certifikátem.
Certifikát ve~formátu PFX je načítán z~konfigurované cesty a~připojen do~HTTP handleru.
QR kód s~přechodovým tokenem, jejž host snímá vlastním zařízením, je vykreslen knihovnou \texttt{QRCoder}.

Hlášení cizinců cizinecké policii \texttt{Ubyport} (časový limit 30~sekund) probíhá nad protokolem SOAP s~HTTP Basic autentizací zajištěnou typem \texttt{NetworkCredential}.
Vzhledem k~tomu, že WSDL definice protistrany prošla v~průběhu vývoje několika nekompatibilními změnami, je SOAP obálka konstruována ručně bez~generovaného klienta.
Toto řešení zachovává plnou kontrolu nad tvarem odesílaných hlaviček a~elementů a~umožňuje reagovat na~změny protistrany bodovou úpravou jediné metody namísto regenerace celého proxy.

\section{Klientské aplikace}
\label{sec:impl-clients}
Všichni tři klienti systému\,--\,webové rozhraní recepce s~veřejným rezervačním portálem, tabletová aplikace pro hosty u~recepčního pultu a~mobilní aplikace pro provozní personál\,--\,jsou realizováni jako jediný projekt nad frameworkem \texttt{Angular}\footnote{Angular -- dostupné z:~\url{https://angular.dev/}.} ve~verzi~21.
Projekt je sestavován do~podoby progresivní webové aplikace a~distribuován z~téhož bundlu prostřednictvím servisního workera.
Cílový klient je rozlišen až na~úrovni směrování: kořenový \texttt{Router} v~\texttt{app.routes.ts} líně načítá tři odlišné podstromy odpovídající jednotlivým klientům\,--\,\texttt{/public}, \texttt{/staff} a~\texttt{/reception-tablet}\,--\,čímž si každý klient po~vstupu na~svou URL stáhne pouze tu část kódu, kterou skutečně potřebuje.

\subsection{Společná infrastruktura}
\label{subsec:impl-clients-shared}
Sdílená vrstva projektu obsahuje typové definice doménových modelů a~tenkého klienta nad třídou \texttt{HttpClient}.
Tento klient je rozšířen o~tři interceptory: první přidává hlavičku \texttt{Authorization} s~přístupovým tokenem, druhý zajišťuje transparentní obnovu vypršelých tokenů přes refresh-token a~třetí překládá serverové chybové odpovědi na~typovaný objekt.

Autentizační služba řeší tři související úlohy.
Perzistence relace probíhá v~úložišti prohlížeče \texttt{localStorage}.
Koordinace obnovy tokenů napříč otevřenými záložkami je realizována nad rozhraními \texttt{BroadcastChannel}\footnote{Broadcast Channel API -- dostupné z:~\url{https://developer.mozilla.org/en-US/docs/Web/API/Broadcast_Channel_API}.} a~\texttt{Web~Locks}\footnote{Web Locks API -- dostupné z:~\url{https://developer.mozilla.org/en-US/docs/Web/API/Web_Locks_API}.}, takže refresh-token není současně použit dvěma instancemi téhož klienta.
Přihlašování probíhá výhradně bezheslově pomocí passkeys.

Stav komponent je důsledně modelován pomocí signálů\,--\,primitiva \texttt{signal}, odvozených hodnot a~asynchronních zdrojů.
Změna detekce probíhá v~bezzónovém režimu konfigurovaném přes \texttt{provideZonelessChangeDetection}, takže běhové prostředí \texttt{zone.js} není v~aplikaci zavedeno vůbec.
Namísto starších reaktivních forem jsou složitější vstupní obrazovky postaveny nad novým signálovým rozhraním \texttt{@angular/forms/signals}, které umožňuje deklarativně skládat typové i~podmíněné validátory přímo nad signálovým modelem formuláře.

Uživatelské rozhraní je sestaveno nad knihovnou \texttt{PrimeNG}\footnote{PrimeNG -- dostupné z:~\url{https://primeng.org/}.} s~tématem \emph{Aura}, ikonografie pochází z~balíčku \texttt{PrimeIcons} a~překlady knihovny jsou inicializovány buď do~češtiny, nebo do~angličtiny podle aktivního \texttt{LOCALE\_ID}.
Lokalizace samotné aplikace je řešena vestavěným nástrojem \texttt{@angular/localize}\footnote{Angular Internationalization -- dostupné z:~\url{https://angular.dev/guide/i18n}.} s~bilingvním pokrytím přihlašovací stránky a~veřejného portálu, zatímco zbytek personální agendy je úmyslně zafixován do~češtiny, neboť cílovou skupinou jsou pouze tuzemští zaměstnanci kempu.

\subsection{Veřejný rezervační portál}
\label{subsec:impl-clients-public}
Veřejný portál vystavuje odděleným vstupním bodem pouze průvodce novou rezervací, detail rezervace a~vícekrokový online check-in.
Průvodce je postaven nad komponentou \texttt{Stepper} a~obsahuje kroky pro výběr období, výběr míst, kontaktní údaje a~rekapitulaci.
Online check-in má validaci odstupňovanou podle státní příslušnosti a~věku jednotlivých hostů.

\subsection{Personální desktopová aplikace}
\label{subsec:impl-clients-desktop}
Personální podstrom \texttt{/staff} je vnitřně rozdělen na~větev \texttt{desktop} a~\texttt{mobile}.
Přepínání mezi nimi je vyhodnoceno na základě \texttt{mediaQuery}.
Rozsah personální agendy v~desktopové variantě sahá od~plánovacího kalendáře a~rezervací přes osmikrokový průvodce sestavením účtu, evidenci faktur a~účetní uzávěrky až po~systémové nastavení, správu uživatelů a~přístupových karet.
Přístup k~jednotlivým trasám je řízen rolovými strážci tras s~výčtem rolí recepční, účetní a~správce.

\subsection{Tabletová aplikace recepce}
\label{subsec:impl-clients-tablet}
Tablet umístěný na~recepčním pultu komunikuje s~odděleným serverem recepce v~reálném čase přes klienta \texttt{socket.io-client}\footnote{Socket.IO Client -- dostupné z:~\url{https://socket.io/docs/v4/client-api/}.}.
Po~spuštění je tablet spárován s~konkrétní recepční stanicí načtením QR kódu, který je v~desktopovém rozhraní recepce vykreslován balíčkem \texttt{qrcode}\footnote{qrcode -- dostupné z:~\url{https://github.com/soldair/node-qrcode}.}.
V~průběhu relace tablet přijímá události s~rekapitulací pobytu a~vrací zpět zachycené digitální podpisy hostů; pořizování podpisu je realizováno vlastní obálkou nad knihovnou \texttt{signature\_pad}\footnote{Signature Pad -- dostupné z:~\url{https://github.com/szimek/signature_pad}.}.
Tablet zpět vrací digitální podpisy cizinců a~ověření občanky přes službu \texttt{eDoklady}.

\subsection{Mobilní aplikace provozního personálu}
\label{subsec:impl-clients-mobile}
Mobilní klient pro provozní personál v~obrazovce skenování využívá výhradně nativní prohlížečová rozhraní \texttt{BarcodeDetector}\footnote{Barcode Detection API -- dostupné z:~\url{https://developer.mozilla.org/en-US/docs/Web/API/BarcodeDetector}.} a~\texttt{TextDetector}\footnote{Text Detection API -- dostupné z:~\url{https://wicg.github.io/shape-detection-api/text.html}.} provozovaná ve \texttt{Web Workerech}\footnote{Web Workers API -- dostupné z:~\url{https://developer.mozilla.org/en-US/docs/Web/API/Web_Workers_API}.}.
Rozpoznávání QR kódů z~ubytovacích míst i~registračních značek vozidel tak probíhá kompletně lokálně v~zařízení.
Obě uvedená rozhraní jsou aktuálně dostupná pouze v~prohlížečích postavených na~enginu Chromium; vzhledem k~unifikovanému parku služebních zařízení (Android s~prohlížečem Chrome) jde o~akceptovatelné omezení.
Rozpoznané hodnoty jsou následně dotazovány proti aplikačnímu rozhraní.

\subsection{Sestavení a~distribuce}
\label{subsec:impl-clients-build}
Ze~sdíleného zdrojového stromu vzniká jediný optimalizovaný webový bundle.
Ve~výrobní konfiguraci je doplněn o~servisního workera generovaného z~konfigurace, čímž se~všichni tři klienti stávají instalovatelnými na~koncových zařízeních jako \texttt{PWA}\footnote{Progresivní webová aplikace.} bez~nutnosti distribuovat samostatné nativní balíčky pro mobilní platformy ani desktopové prostředí.
Pro anglickou jazykovou variantu se~tentýž zdrojový kód sestavuje s~odlišným \texttt{baseHref} a~předpřeloženými řetězci do~paralelního výstupu; z~hlediska procesu sestavení i~doručení jde o~jedinou systematickou odlišnost mezi klienty.

\section{Ovládání vjezdových závor}
\label{sec:impl-gate}
Vjezdová služba je realizována jako aplikace v~jazyce \texttt{C\#} nad platformou \texttt{.NET~10}, hostovaná pomocí balíčku \texttt{Microsoft.Extensions.Hosting.WindowsServices} jako nativní Windows služba spouštěná z~\emph{Service Control Manageru}.
REST rozhraní pro příjem požadavků z~backendu je vystaveno prostřednictvím \emph{ASP.NET~Core~Minimal~API} nad webovým serverem \texttt{Kestrel}.
Klientům jsou vystaveny pouze dva idempotentní endpointy nad cestou \texttt{/api/v1/cards/\{key\}}: metoda \texttt{PUT} provádí vložení nebo aktualizaci RFID karty a~metoda \texttt{DELETE} její odstranění.
Idempotence obou operací umožňuje backendu opakovat selhané volání, aniž by hrozila desynchronizace stavu mezi systémy.

Pro integraci s~legacy systémem Kartoun, jehož datovou vrstvu tvoří dvojice zaheslovaných souborů formátu \texttt{Microsoft Access .mdb} (databáze uživatelů a~databáze auditních událostí), je využita knihovna \texttt{Microsoft.Jet.OLEDB.4.0}\footnote{Microsoft Jet Database Engine -- dostupné z:~\url{https://learn.microsoft.com/en-us/office/client-developer/access/desktop-database-reference/microsoft-ole-db-provider-for-microsoft-jet}.}.
Modernější poskytovatel \emph{ACE~OLEDB} zde nepřichází v~úvahu: existující soubory pocházejí z~období Access~97--2003 a~jejich heslovací schéma zachované ovladačem Jet není s~poskytovatelem ACE plně kompatibilní.
Veškeré dotazy jsou formulovány jako parametrizované příkazy \texttt{OleDbCommand}, čímž je eliminováno riziko útoku typu SQL injection.
Vzhledem k~tomu, že jmenovaný OLE~DB poskytovatel je dostupný výhradně v~32bitové variantě, je celá aplikace kompilována pro architekturu \texttt{win-x86}.

Časové sloupce původního schématu (\texttt{ValidFrom}, \texttt{ValidTo}, \texttt{DateTime}) jsou ukládány jako celá čísla představující počet sekund od~proprietární epochy 1.\,1.\,2000; konverze je zapouzdřena v~pomocném typu \texttt{LegacyTime}.
Každá úspěšná modifikace tabulky \texttt{Card} je doplněna o~odpovídající záznam v~tabulce \texttt{Events}, čímž je zachována konzistence s~auditním modelem původní aplikace.

Při souběžném zápisu uzamyká Jet celý soubor a~konkurenční operace selže okamžitě.
Toto chování je obslouženo retry strategií s~exponenciálním zpožděním omezeným na~několik desítek pokusů, čímž se~překlene typicky milisekundové okno zámku držené jiným procesem.
Strukturované logování je opět zajištěno knihovnou \texttt{Serilog} s~výstupem do~lokálního souborového sinku rotovaného po~dnech.
V~případě nevratného selhání procesu zajišťuje jeho restart \emph{Service Control Manager} podle nakonfigurované obnovovací politiky.

\section{Lokální tiskový server}
\label{sec:impl-print}
Tisková služba je realizována stejným způsobem jako ovladač vjezdových závor, viz \ref{sec:impl-gate}.
Webový server \texttt{Kestrel} naslouchá pouze na~loopback rozhraní a~přijímá požadavky pouze z~webové aplikace povolené zásadami \texttt{CORS}.
Vazba na~loopback eliminuje vystavení tiskového API mimo lokální stanici, a~je proto základním obranným prvkem této služby.

Rozhraní vystavuje dva koncové body: výpis nainstalovaných tiskáren získaný přes \texttt{System.Drawing.Printing.PrinterSettings} a~odeslání PDF dokumentu na~zvolenou tiskárnu.
Samotný tisk je delegován na~externí nástroj \texttt{SumatraPDF}\footnote{SumatraPDF -- dostupné z:~\url{https://www.sumatrapdfreader.org/}.} spouštěný v~tichém režimu jako podproces, čímž se~aplikace vyhýbá přímé integraci s~tiskovými ovladači systému Windows a~deleguje renderování PDF na~ověřenou implementaci.
Selhání podprocesu (nenulový návratový kód, překročení časového limitu) jsou propagována zpět klientovi jako odpověď stavového kódu 502 a~zároveň zaznamenána do~logu prostřednictvím knihovny \texttt{Serilog}.
